\chapter{Kalman Filter}\label{chap:kalman}

\section{Introduction}

The Kalman filter estimates the state of a linear dynamical system from a
sequence of noisy measurements. In the robotics and signal-processing
literature the state at time $t$ is usually written $\mathbf{x}_t$ and
the measurement $\mathbf{z}_t$ (see, e.g., \cite{thrun2005probabilistic}).
The Deep Kalman Filter paper \cite{krishnan2015deep} follows the VAE
convention instead, and I do the same throughout this report:
$\mathbf{z}_t$ is the \emph{latent} state and $\mathbf{x}_t$ the
\emph{observed} measurement.

\begin{story}[Control input $\mathbf{u}_{t-1}$ versus $\mathbf{u}_{t}$]
Krishnan et al.\ write the control input that drives the transition from
$t-1$ to $t$ as $\mathbf{u}_{t-1}$, but when they define their
variational models, $\mathbf{u}_{t}$ is synchronous with
$\mathbf{z}_{t}$. I keep the convention of the Kalman-filter literature
and write $\mathbf{u}_{t}$ for the input acting on the transition into
$\mathbf{z}_{t}$. This is only a matter of indexing.
\end{story}

\section{System Model}

The system is described by an initial belief, a state-transition
equation and an observation equation.

\paragraph{Initial state.}
$\mathbf{z}_{0} \sim \mathcal{N}(\boldsymbol{\mu}_{0},
\boldsymbol{\Sigma}_{0})$.

\paragraph{State transition.} For $t \geq 1$,
\begin{equation*}
    \mathbf{z}_t = \mathbf{G}_t \mathbf{z}_{t-1} + \mathbf{B}_t
    \mathbf{u}_t + \boldsymbol{\varepsilon}_t,
    \qquad
    \boldsymbol{\varepsilon}_t \sim \mathcal{N}(\mathbf{0},
    \boldsymbol{\Sigma}_{t}),
\end{equation*}
where $\mathbf{G}_t$ is the state-transition matrix, $\mathbf{B}_t$ the
control-input matrix, $\mathbf{u}_t$ the (known) control input and
$\boldsymbol{\varepsilon}_t$ the process noise. Equivalently,
\begin{equation}
    p(\mathbf{z}_t \mid \mathbf{z}_{t-1}, \mathbf{u}_t) =
    \mathcal{N}(\mathbf{z}_t;\ \mathbf{G}_t \mathbf{z}_{t-1} +
    \mathbf{B}_t \mathbf{u}_t,\ \boldsymbol{\Sigma}_{t}).
    \label{eq:1}
\end{equation}

\paragraph{Observation.} For $t \geq 1$,
\begin{equation*}
    \mathbf{x}_t = \mathbf{F}_t \mathbf{z}_t + \boldsymbol{\eta}_t,
    \qquad
    \boldsymbol{\eta}_t \sim \mathcal{N}(\mathbf{0},
    \boldsymbol{\Gamma}_{t}),
\end{equation*}
where $\mathbf{F}_t$ is the observation matrix and
$\boldsymbol{\eta}_t$ the measurement noise. Equivalently,
\begin{equation}
    p(\mathbf{x}_t \mid \mathbf{z}_t) =
    \mathcal{N}(\mathbf{x}_t;\ \mathbf{F}_t \mathbf{z}_{t},\
    \boldsymbol{\Gamma}_{t}).
    \label{eq:2}
\end{equation}

\section{Kalman Filter Algorithm}

Write $\hat{\mathbf{z}}_{t}$ and $\mathbf{P}_{t}$ for the mean and
covariance of the state given the measurements up to time $t$. The filter
alternates two steps.

\subsection{Prediction Step}
\begin{itemize}
    \item \textbf{Predicted (a priori) mean:}\quad
    $\bar{\mathbf{z}}_{t} = \mathbf{G}_t \hat{\mathbf{z}}_{t-1} +
    \mathbf{B}_t \mathbf{u}_t$.
    \item \textbf{Predicted (a priori) covariance:}\quad
    $\bar{\mathbf{P}}_{t} = \mathbf{G}_t \mathbf{P}_{t-1}
    \mathbf{G}_t^\top + \boldsymbol{\Sigma}_{t}$.
\end{itemize}
These are the mean and covariance of the state at time $t$ given the
measurements up to time $t-1$.

\subsection{Update Step}
\begin{itemize}
    \item \textbf{Innovation (measurement residual):}\quad
    $\mathbf{y}_t = \mathbf{x}_t - \mathbf{F}_t \bar{\mathbf{z}}_{t}$.
    \item \textbf{Innovation covariance:}\quad
    $\mathbf{S}_t = \mathbf{F}_t \bar{\mathbf{P}}_{t} \mathbf{F}_t^\top
    + \boldsymbol{\Gamma}_{t}$.
    \item \textbf{Kalman gain:}\quad
    $\mathbf{K}_t = \bar{\mathbf{P}}_{t} \mathbf{F}_t^\top
    \mathbf{S}_t^{-1}$.
    \item \textbf{Updated (a posteriori) mean:}\quad
    $\hat{\mathbf{z}}_{t} = \bar{\mathbf{z}}_{t} + \mathbf{K}_t
    \mathbf{y}_t$.
    \item \textbf{Updated (a posteriori) covariance:}\quad
    $\mathbf{P}_{t} = (\mathbf{I} - \mathbf{K}_t \mathbf{F}_t)
    \bar{\mathbf{P}}_{t}$.
\end{itemize}

\section{The probabilistic reading of the Kalman filter}

Because every distribution in the model is Gaussian and every map is
linear, all posterior distributions are Gaussian and can be computed in
closed form. The two steps above compute exactly
\begin{equation*}
    p(\mathbf{z}_t \mid \mathbf{x}_{1:t-1}, \mathbf{u}_{1:t}) =
    \mathcal{N}(\bar{\mathbf{z}}_{t}, \bar{\mathbf{P}}_{t})
    \qquad\text{and}\qquad
    p(\mathbf{z}_t \mid \mathbf{x}_{1:t}, \mathbf{u}_{1:t}) =
    \mathcal{N}(\hat{\mathbf{z}}_{t}, \mathbf{P}_{t}).
\end{equation*}
The innovations also give the exact log-likelihood of the measurements,
\begin{equation*}
    \log p(\mathbf{x}_{1:T} \mid \mathbf{u}_{1:T}) =
    \sum_{t=1}^{T} \log \mathcal{N}(\mathbf{y}_t;\ \mathbf{0},\
    \mathbf{S}_t),
\end{equation*}
and a backward pass (the Rauch--Tung--Striebel smoother) gives the
\emph{smoothing} distributions
$p(\mathbf{z}_t \mid \mathbf{x}_{1:T}, \mathbf{u}_{1:T})$, which use the
future measurements as well.

\begin{inwords}
In the linear-Gaussian case, \enquote{filtering} and \enquote{smoothing}
are exact Bayesian inference, done with matrix algebra. The Deep Kalman
Filter keeps the same probabilistic structure but makes the transition
and emission non-linear. Then none of these distributions has a closed
form any more. The DKF replaces exact inference by a neural network that
\emph{approximates} the smoothing posterior, and replaces the exact
log-likelihood by a lower bound on it. Chapters~\ref{chap:vae} and
\ref{chap:dkf} explain how.
\end{inwords}

\section{Summary}

The Kalman filter alternates a prediction step, which propagates the
belief through the dynamics, and an update step, which corrects it with
the new measurement. For linear-Gaussian systems this is exact inference,
and the resulting estimate is optimal in the mean-squared-error sense.
The Deep Kalman Filter replaces all the linear maps by non-linear maps
parameterised by neural networks, at the price of approximate inference.
